\section{Related Work}
\label{sec:related}

\paragraph{Audio-driven Single Image Talking Face Generation.} Early works~\cite{lipgan, wav2lip, videoretalking} mainly focus on producing accurate lip motion with a perception discriminator. Since the real videos contain many different motions, ATVGnet~\cite{atvgnet} uses the facial landmark as the intermediate representation to generate the video frames. A similar approach has been proposed by MakeItTalk~\cite{zhou2020makelttalk}, differently, it disentangles the content and speaker information from the input audio signal. 
Since facial landmarks are still a highly coupled space, generating the talking head in the disentangled space is also popular recently. PC-AVS~\cite{pcavs} disentangles the head pose and expression using implicit latent code. However, it can only produce low-resolution image and need the control signal from another video.
Audio2Head~\cite{wang2021audio2head} and Wang~\etal~\cite{wang2021one} get inspiration from the video-driven method~\cite{fomm} to produce the talking-head face. However, these head movements are still not vivid and produce distorted faces with inaccurate identities. Although there are some previous works~\cite{hdtf, ren2021pirenderer} use 3DMMs as an intermediate representation, their method still faces the problem of inaccurate expressions~\cite{ren2021pirenderer} and obvious artifacts~\cite{hdtf}.
 
\paragraph{Audio-driven Video Portrait.} Our task is also related to visual dubbing, which aims to edit a portrait video through audio. Different from audio-driven single image talking face generation, this task is typically required to be trained and edited on the specific video. Following previous work of deep video portrait~\cite{dvp}, these methods utilize 3DMM information for face reconstruction and animation. AudioDVP~\cite{audiodvp}, NVP~\cite{nvp}, AD-NeRF~\cite{adnerf} learn to reenact the expression to edit the mouth shape. Beyond lip movement, \ie, the head motions~\cite{facial, lsp}, emotional talking face~\cite{evp} also get attention. The 3DMM-based method plays an important role in these tasks since it is practical to fit the 3DMM parameters from a video clip. Although these methods achieve satisfactory results in personalized video, their method can not be applied to arbitrary photos and in the wild audio.

 \paragraph{Video-Driven Single Image Talking Face Generation.} This task is also known as face reenactment or face animation, which aims to transfer the motion of the source image to the target person. It has been widely explored~\cite{fomm, mraa, ren2021pirenderer, facevid2vid, fewshotvid2vid, styleheat, dagan, tps, lia, dpe} recently. Previous works also learn a shared intermediate motion representation from the source image and the target, which can be roughly divided into the landmark~\cite{fewshotvid2vid} and the unsupervised landmark-based methods~\cite{fomm, facevid2vid,dagan,tps}, 3DMM based methods~\cite{ren2021pirenderer, styleheat, doukas2020headgan} and the latent animation~\cite{lia, mallya2022implicit}. This task is much easier than our task since it contains the motion in the same domain. Our face render is also inspired by the method of unsupervised landmark-based method~\cite{facevid2vid} and 3DMM-based method~\cite{ren2021pirenderer} to map the learned coefficient to generate the real video. However, they are not focused on generating realistic motion coefficients.